\section{Introduction}
to fix
External or historical cohorts are increasingly used to augment small
concurrent studies in medicine.
When the external and concurrent sources are sufficiently similar, borrowing
can improve efficiency and stabilize estimation.
When the sources differ in meaningful ways, however, naive borrowing can
introduce bias and distort uncertainty quantification
\citep{viele2014use}.
This tension is especially important in Alzheimer's disease research, where
prospective follow-up is expensive, clinically meaningful outcomes accrue
slowly, and related cohorts often differ in enrollment criteria, disease
severity, and assessment practice.

Bayesian dynamic borrowing (BDB) provides a principled framework for
adaptively controlling how much external information should influence
current-study inference
\citep{ibrahim2000power,hobbs2011commensurate,neuenschwander2010summarizing,
schmidli2014robust}.
Classical BDB methods such as the power prior, commensurate prior, and robust
MAP prior have strong statistical motivation, but in practice they still
require repeated model fitting for each new dataset pair and often summarize
cross-source compatibility through one or a few low-dimensional quantities.
That can be restrictive in settings where source mismatch is multi-faceted.

Amortized neural posterior estimation (NPE) offers a different route.
Instead of rerunning MCMC or numerical optimization for every new
current/external pair, the idea is to train a neural network offline on many
simulated borrowing problems and then use the trained network to output an
approximate posterior in a single forward pass
\citep{papamakarios2019sequential,greenberg2019automatic,cranmer2020frontier,
lueckmann2021benchmarking}.
The key design question is how to represent the two sources so that the
network sees both the target estimand and the degree of non-exchangeability.

In this paper we focus on a summary-based NPE approach for dynamic borrowing.
We first study the method in a controlled simulation setting spanning six
exchangeability regimes, including covariate shift, outcome drift, and joint
non-exchangeability.
We then examine a real-data borrowing case study using the Alzheimer's Disease
Neuroimaging Initiative (ADNI) \citep{petersen2010alzheimer,weiner2017alzheimer}.
Specifically, ADNI1 is treated as an external cohort and ADNIGO/ADNI2 as a
concurrent cohort among baseline mild cognitive impairment (MCI) subjects, with
24-month conversion to dementia as the outcome.
This real-world application is not presented as a randomized treatment-effect
analysis; rather, it is used to study borrowing under cohort shift in a
clinically meaningful longitudinal dataset.

The ADNI application serves two purposes.
First, it provides an empirical check that the borrowing problem is real in
practice: the external and concurrent cohorts have different event rates and
therefore are not directly exchangeable.
Second, it gives a setting in which we can adapt the NPE framework to infer a
concurrent risk-level parameter while allowing the external cohort to differ
through a nuisance shift parameter.
This keeps the spirit of dynamic borrowing while avoiding an overstatement
that the real-data example exactly matches the treatment-effect simulations.

The remainder of the paper is organized as follows.
Section~\ref{sec:related} reviews dynamic borrowing, balancing methods, and
simulation-based inference.
Section~\ref{sec:problem} introduces the simulation regimes and the ADNI
borrowing setup.
Section~\ref{sec:figures_plan} presents the empirical design, summarizes the
current simulation workflow, and reports the ADNI real-world application.

% -------------------------------------------------------
% REQUIRED MLHC SECTION
% -------------------------------------------------------
\subsection*{Generalizable Insights about Machine Learning in the Context
             of Healthcare}

\begin{itemize}[leftmargin=1.5em, itemsep=2pt]
  \item \textbf{Amortized inference can make adaptive borrowing practical in
        healthcare studies with repeated external comparisons.}
        Training an NPE model offline allows approximate posterior inference
        for new cohort pairs without rerunning expensive Bayesian computation
        from scratch each time.

  \item \textbf{Real-world borrowing problems often appear first as calibration
        mismatch rather than discrimination failure.}
        In the ADNI application, external-only models rank subjects well but
        systematically overestimate risk in the concurrent cohort, showing that
        borrowing decisions should account for cohort shift even when predictive
        AUC remains high.

  \item \textbf{A combined simulation-plus-case-study design is useful for
        evaluating healthcare borrowing methods.}
        Controlled simulations make it possible to assess bias and error under
        known exchangeability violations, while a real-world case study checks
        whether the same issues appear in practice and whether the method
        behaves sensibly on clinically meaningful data.
\end{itemize}
